% Generated from current MATLAB exports. Include in the correct main/SI sequence.
\begingroup\small\setlength{\tabcolsep}{3pt}
\begin{longtable}{>{\raggedright\arraybackslash}p{0.1700\linewidth}>{\raggedright\arraybackslash}p{0.1650\linewidth}>{\raggedright\arraybackslash}p{0.1650\linewidth}>{\raggedright\arraybackslash}p{0.1650\linewidth}>{\raggedleft\arraybackslash}p{0.1650\linewidth}}
\caption{Completed retrospective peak origins, training/test lengths, and original random seeds.} \label{tab:oos-splits} \\
\toprule
Country & Origin & Cutoff date & Train/test days & Seed \\
\midrule
\endfirsthead
\caption[]{Completed retrospective peak origins, training/test lengths, and original random seeds.} \\
\toprule
Country & Origin & Cutoff date & Train/test days & Seed \\
\midrule
\endhead
\midrule
\multicolumn{5}{r}{Continued on next page} \\
\midrule
\endfoot
\bottomrule
\endlastfoot
Korea & Second peak (Wave 1) & 2020-06-24 & 130/48 & 1 \\
Japan & Observed peak (Wave 1) & 2020-05-04 & 56/55 & 123 \\
Pakistan & Observed peak (Wave 1) & 2020-07-02 & 108/103 & 1 \\
\end{longtable}
\par\noindent\footnotesize\textit{Note.} All three completed original peak forecasts are retained. Pakistan has a budget-limited initializer; Korea/Japan forms the separate convergence subset. No forecast row is selected by its held-out error. Intervals are latent ensemble quantiles; the residual score is descriptive.
\endgroup
